\documentclass[../my_ntust_thesis.tex]{subfiles}
\begin{document}
\chapter{Proposed Method}
\label{sec:method}

\textcolor{blue}{This section details the proposed dynamic timing adjustment algorithm. To resolve the unstable latency during packet transmission between the PNF and VNF, we propose a Dual-Loop Timing Control Architecture.}

\section{Dual-Loop Architecture}
\label{subsec:dual_loop_architecture}

\textcolor{blue}{The timing control of this system adopts a ``policy-execution separation'' dual-loop architecture, ensuring that the timing advancement on the VNF side can flexibly cope with network jitter without inducing unexpected oscillations:
\begin{enumerate}
    \item \textbf{Convergence Optimization Loop (Policy):} Dynamically analyzes the statistical data reported by the PNF through the convergence optimization algorithm to determine the global target slot-ahead $s_{\text{ahead}}$.
    \item \textbf{Timing Actuator Loop (Execution):} Executed by an independent VNF timing thread. This thread is the only unit in the system capable of altering the absolute sleep time and slot propagation, translating the upper-layer policy target and low-layer microsecond phase deviations into actual timing advancement.
\end{enumerate}}

\section{VNF Timing Actuator}
\label{subsec:vnf_timing_actuator}

\textcolor{blue}{The core function of the \textbf{VNF Timing Actuator} is isolation and precise phase execution. Initially, we define the basic mathematical parameters used in the system. Let $\mu$ denote the subcarrier spacing parameter (Numerology), and the corresponding slot duration (in microseconds) is defined as $T_{\text{slot}}$:
\begin{equation} \label{eq:tslot}
T_{\text{slot}} = \frac{1000}{2^\mu} 
\end{equation}
Let $s_{\text{ahead}}$ be the target slot-ahead value, and $T_{\text{adv}}$ be the accumulated total phase advance in microseconds. For microsecond-level phase correction, let $\Delta_{\text{slot}}$ be the integer slot correction provided by the lower layer, and $\Delta_{\text{pending}}$ be the sub-slot phase deviation (in \textmu s). The algorithm calculates the next absolute wake-up time $T_{\text{next}}$:
\begin{equation} \label{eq:tnext}
T_{\text{next}}^{(k)} = T_{\text{next}}^{(k-1)} + T_{\text{slot}} + \Delta_{\text{pending}}
\end{equation}
Before receiving the initial timing anchor from the PNF, the thread blocks execution via a condition variable. Once active, the thread operates within a mutex-locked environment and triggers sleep according to the absolute monotonic clock calculated in \eqref{eq:tnext}, mitigating long-term error accumulation. The condensed control logic is presented in Algorithm~\ref{alg:vnf_timing_thread}.}

\begin{algorithm}[H]
\caption{VNF Timing Actuator Thread}
\label{alg:vnf_timing_thread}
\begin{algorithmic}[1]
\Procedure{TimingActuator}{$s_{\text{ahead}}, \mu$}
  \State Wait for \texttt{timing\_info} condition variable
  \State $T_{\text{next}} \gets \text{Now}()$
  \State $ID_{\text{DEC}} \gets \mathrm{DEC}(\mu, sfn, slot)$
  \While{$running$}
    \State \textbf{Lock} Mutex
    \State $T_{\text{adv}} \gets s_{\text{ahead}} \cdot T_{\text{slot}}$
    \If{$\Delta_{\text{slot}} \neq 0$} \Comment{Integer correction}
      \State $ID_{\text{DEC}} \mathrel{+}= \Delta_{\text{slot}}$
      \If{$sync\_locked$} \State $T_{\text{adv}} \mathrel{+}= \Delta_{\text{slot}}$ \EndIf
      \State $\Delta_{\text{slot}} \gets 0$
    \EndIf
    \State $\Delta_{\text{current}} \gets \Delta_{\text{pending}}$ \Comment{Sub-slot correction}
    \If{$sync\_locked$} \State $T_{\text{adv}} \mathrel{-}= \Delta_{\text{current}}$ \EndIf
    \State $\Delta_{\text{pending}} \gets 0$
    \State \textbf{Unlock} Mutex

    \State $T_{\text{next}} \mathrel{+}= T_{\text{slot}} + \Delta_{\text{current}}$
    \State \Call{SleepUntil}{$T_{\text{next}}$}
    
    \State \Call{HandlePeriodicSync}{}
    \State $target \gets \mathrm{wrap}(ID_{\text{DEC}} + s_{\text{ahead}})$
    \While{$last\_mac \neq target$} \Comment{Indication catch-up}
      \State $last\_mac \mathrel{+}= 1$ (wrapped)
      \State \Call{DeliverToMAC}{$last\_mac$}
    \EndWhile
    \State $ID_{\text{DEC}} \mathrel{+}= 1$ (wrapped)
  \EndWhile
\EndProcedure
\end{algorithmic}
\end{algorithm}

\section{Convergence Optimization Policy}
\label{subsec:convergence_optimization}

\textcolor{blue}{To counteract fluctuating network latency and maintain scheduling robustness, this paper proposes a dynamic convergence optimization policy. Upon receiving the timing statistics reported by the PNF, the system incorporates an Exponentially Weighted Moving Average (EWMA) model to estimate the mean latency $\bar{L}$ and jitter variance $V_{\text{jitter}}$. Given the worst-case late arrival $L_{\text{worst}}$ in the current observation period, the iterative update formulas are:
\begin{align}
\label{eq:ewma_mean}
\bar{L}^{(k)} ={}& \bar{L}^{(k-1)} + \frac{L_{\text{worst}} - \bar{L}^{(k-1)}}{8} \\
\label{eq:ewma_var}
V_{\text{jitter}}^{(k)} ={}& V_{\text{jitter}}^{(k-1)} + \frac{|L_{\text{worst}} - \bar{L}^{(k-1)}| - V_{\text{jitter}}^{(k-1)}}{4}
\end{align}}

\textcolor{blue}{Before making tracking adjustments, a dynamic safety margin $M_{\text{safe}}$ is established to cover the networking jitter, and an absolute boundary $B_{\text{abs}}$ is enforced to limit overly aggressive delays:
\begin{align}
M_{\text{safe}} ={}& \max\left(1.5 \cdot T_{\text{slot}}, 1.5 \cdot V_{\text{jitter}}\right) \\
B_{\text{abs}} ={}& 2 \cdot T_{\text{slot}}
\end{align}
Furthermore, the algorithm designs multiple \textbf{Panic Criteria} triggers (including deadline edge $panic_{\text{edge}}$, severe jitter $panic_{\text{jitter}}$, statistical anomaly $anomaly_{\text{stat}}$, and strict violation $violation_{\text{strict}}$). To prevent unexpected oscillation during target transitions, a \textbf{Directionality Guard} is introduced. Additionally, a \textbf{Penalty Counter \& Anti-bounce} mechanism filters out pseudo-stable low-latency states to prevent prematurely lowering $s_{\text{ahead}}$. When panic events are detected, the algorithm initiates a fast-attack mode, radically increasing $S_{\text{target}}$ to ensure connection continuity. Conversely, when redundant phase advance exists continuously without penalties, a slow-decay backoff is triggered. The overall dynamic adjustment is condensed in Algorithm~\ref{alg:convergence_opt}.}


\begin{algorithm}
\caption{Delay Manager: Dynamic Scheduler Trigger Timing}
\label{alg:timing_optimization}
\begin{algorithmic}[1]
\Require 
    \Statex $t_{late}$: Worst late timing from PHY (Input)
    \Statex $T_{slot}$: Slot duration in microseconds
    \Statex $S_{curr}$: Current slots ahead (Environment state)
\Ensure 
    \Statex $S_{next}$: Adjusted slots ahead (Output: Higher = Earlier, Lower = Later)

\State \textbf{Update} EWMA Mean Late ($E$) and Jitter Variance ($V$) using $t_{late}$
\State $B_{safe} \gets (S_{curr} - 1) \times T_{slot}$ \Comment{Absolute safe boundary}
\State $C_{panic} \gets (t_{late} \ge T_{slot}) \lor (V > \text{SafeMargin}) \lor (\text{Statistical Anomaly})$

\If{$C_{panic}$}
    \State $S_{next} \gets S_{curr} + \text{StepUp}$ \Comment{\textbf{Earlier:} Fast attack due to violation/jitter}
    \State Reset safety counters and apply reduction penalty
\ElsIf{$t_{late} < -B_{safe}$ \textbf{and} \text{Reduction is not locked}}
    \State $Count_{safe} \gets Count_{safe} + 1$
    \If{$Count_{safe} > \text{Decay Threshold}$}
        \State $S_{next} \gets S_{curr} - 1$ \Comment{\textbf{Later:} Proactive safe latency reduction}
        \State $Count_{safe} \gets 0$
    \EndIf
\Else
    \State $S_{next} \gets S_{curr}$ \Comment{\textbf{Maintain:} Timing is stable within bounds}
\EndIf

\State $S_{next} \gets \max(1, \min(S_{next}, 8))$ \Comment{Apply absolute system bounds}
\State \Return $S_{next}$
\end{algorithmic}
\end{algorithm}

\begin{algorithm}[H]
\caption{Dynamic Convergence Optimization Policy}
\label{alg:convergence_opt}
\begin{algorithmic}[1]
\Procedure{ConvergenceOpt}{$L_{\text{worst}}, T_{\text{slot}}$}
  \State $diff \gets L_{\text{worst}} - \bar{L}$
  \State $\bar{L} \gets \bar{L} + diff / 8$
  \State $V_{\text{jitter}} \gets V_{\text{jitter}} + (|diff| - V_{\text{jitter}}) / 4$
  
  \State $B_{\text{abs}} \gets 2 \cdot T_{\text{slot}}$, $S_{\text{target}} \gets s_{\text{ahead}}$
  
  \State Evaluate logic flags: $panic$, $anomaly_{\text{stat}}$
  \If{$L_{\text{worst}} < 0$} \State $anomaly_{\text{stat}} \gets false$ \EndIf
  
  \State $\Delta T_{\text{adv}} \gets T_{\text{adv}} - last\_T_{\text{adv}}$
  \State $dir\_helpful \gets \mathrm{SignMatch}(\Delta T_{\text{adv}}, L_{\text{worst}})$

  \State $S_{\text{prop}} \gets \mathrm{clamp}(\mathrm{round}(T_{\text{adv}} / T_{\text{slot}}), 1, S_{\max})$
  \If{$L_{\text{worst}} \ge T_{\text{slot}} / 2$} \State $S_{\text{prop}} \mathrel{+}= 1$ \EndIf

  \If{$S_{\text{prop}} \neq S_{\text{target}}$}
    \If{$|S_{\text{prop}} - S_{\text{target}}| > 1$ \textbf{and not} $dir\_helpful$}
      \State $S_{\text{target}} \mathrel{+}= \mathrm{sign}(S_{\text{prop}} - S_{\text{target}})$ 
    \Else
      \State $S_{\text{target}} \gets S_{\text{prop}}$
    \EndIf
  \EndIf

  \If{$panic$ \textbf{or} $anomaly_{\text{stat}}$} \Comment{Fast-Attack}
    \State $step \gets (violation_{\text{strict}}) ? \lfloor L_{\text{worst}} / T_{\text{slot}} \rfloor + 2 : 2$
    \If{$violation_{\text{strict}}$} \State $penalty \mathrel{+}= 10000$ \EndIf
    \State $S_{\text{target}} \gets \min(S_{\text{target}} + step, S_{\max})$
    \State $safe\_cnt \gets 0$
  \ElsIf{$L_{\text{worst}} < -B_{\text{abs}}$} \Comment{Slow-Decay}
    \State $safe\_cnt \mathrel{+}= 1$
    \State $penalty \gets \max(0, penalty - 1)$
    \If{$safe\_cnt > 2000 + penalty$} 
      \State \text{Anti-bounce top-tier hold logic once}
      \State $S_{\text{target}} \mathrel{-}= 1$
      \State $safe\_cnt \gets 0$
    \EndIf
  \Else
    \State Reset counters and top-hold state
  \EndIf

  \State $S_{\text{target}} \gets \mathrm{clamp}(S_{\text{target}}, 1, S_{\max})$
  \If{$S_{\text{target}} \neq s_{\text{ahead}}$} \Comment{Phase drift compensation}
    \State $\bar{L} \mathrel{+}= (S_{\text{target}} - s_{\text{ahead}}) \cdot T_{\text{slot}}$ 
    \State $last\_T_{\text{adv}} \gets T_{\text{adv}}$, $s_{\text{ahead}} \gets S_{\text{target}}$
  \EndIf
\EndProcedure
\end{algorithmic}
\end{algorithm}

\end{document}